\subsection{Euler's Number} 
What is this number $e$? What's it doing in that formula? How does the interest rate get into the exponent?
\begin{itemize}
\item $e$ is an \makeblank[1in] number.
\item Like another irrational number, $\pi$, $e$ is defined to be the number that works in certain formulas.
\item $\pi$ is defined to be the number that works in the \makeblank[1.5in] formula, $C=2\pi r$
\item $e$ is defined to be the number that works in the \makeblank formula.
\end{itemize}
Shouldn't that depend on $r$? How do we get $r$ into the exponent?
\begin{itemize}
\item Start with the compound interest formula, $P(t)=P_0\left(1+\frac{r}{n}\right)^{nt}$
\item The $P_0$ part will stay the same---we'll look at just the $\left(1+\frac{r}{n}\right)^{nt}$ part.
\item Introduce a new variable, $x=\makeblanksmall$
\item Then $\frac{r}{n}=\makeblanksmall$ and $n=\makeblanksmall$
\item So $\left(1+\frac{r}{n}\right)^{nt}=\makeblank[1in]=\makeblank[1in]$
\item If $n$ is \makeblank[1in], $x$ is \makeblank[1in].
\item What happens to \makeblank when $x$ is large?
\end{itemize}
\[
\begin{array}{c|c}
     x&\left(1+\frac{1}{x}\right)^x\\\hline
    10&\htstrut\makeblanksmall\\
   100&\htstrut\makeblanksmall\\
  1000&\htstrut\makeblanksmall\\
 10000&\htstrut\makeblanksmall\\
100000&\htstrut\makeblanksmall
\end{array}
\]
\begin{itemize}
\item $e$ is defined to be the number this gets close to.
\end{itemize}